\SourceLeaf{070}
\[
x=g(\cos G+\sin G\, .\,\BareRoot{-1}),\quad x=g(\cos G-\sin G\, .\,\BareRoot{-1})
\]
valorem 0 obtinet, et proin indefinite per
\[
x-g(\cos G+\sin G\, .\,\BareRoot{-1}),\quad\text{nec non per}\quad x-g(\cos G-\sin G\, .\,\BareRoot{-1})
\]
divisibilis erit. Quoties non est $\sin G=0$, neque $g=0$, hi divisores sunt
inaequales, et proin $X$ etiam per illorum productum
\[
xx-2g\cos G\, .\,x+gg
\]
divisibilis erit, quoties autem vel $\sin G=0$ adeoque $\cos G=\pm1$, vel $g=0$,
illi factores sunt identici scilicet $=x\mp g$. Certum itaque est, functionem $X$
involvere divisorem realem secundi vel primi ordinis, et quum eadem conclusio
rursus de quotiente valeat, $X$ in tales factores complete resolubilis erit. Q.~E.~D.

\ThirdArticle{4}
Quamquam in praecedentibus negotio quod propositum erat, iam plene perfuncti simus, tamen haud superfluum erit, adhuc quaedam de ratiocinatione art. 2
adiicere. A suppositione, $t$ et $u$ pro nullis valoribus indeterminatarum $r,\varphi$ intra limites illic assignatos simul evanescere, ad contradictionem inevitabilem delapsi sumus, unde ipsius suppositionis falsitatem conclusimus. Haec igitur contradictio cessare debet, si revera adsunt valores ipsarum $r,\varphi$, pro quibus $t$ et $u$
simul fiunt $=0$. Quod ut magis illustretur, observamus, pro talibus valoribus
fieri $tt+uu=0$, adeoque ipsam $y$ infinitam, unde haud amplius licebit, integrale duplex $\iint y\,\mathrm dr\,\mathrm d\varphi$ tamquam quantitatem assignabilem tractare. Generaliter quidem loquendo, denotantibus $\xi,\eta,\zeta$ indefinite coordinatas punctorum in
spatio, integrale $\iint y\,\mathrm dr\,\mathrm d\varphi$ exhibet volumen solidi, quod continetur inter quinque plana, quorum aequationes sunt
\[
\xi=0,\quad\eta=0,\quad\zeta=0,\quad\xi=R,\quad\eta=360^\circ
\]
atque superficiem, cuius aequatio $\zeta=y$, considerando eas partes tamquam negativas, in quibus coordinatae $\zeta$ sunt negativae. Sed tacite hic subintelligitur,
superficiem sextam esse \textit{continuum}, qua conditione cessante, dum $y$ evadit infinita, utique fieri potest, ut conceptus ille sensu careat. In tali casu de integrali
$\iint y\,\mathrm dr\,\mathrm d\varphi$ colligendo sermo esse nequit, neque adeo mirandum est, operationes
analyticas coeco calculo ad inania applicatas ad absurda perducere.

\SourceLeaf{071}
Integratio $\int y\,\mathrm d\varphi=\frac{tu'-ut'}{r(tt+uu)}$ eatenus tantum est integratio vera, i. e. summatio, quatenus inter limites, per quos extenditur, $y$ ubique est quantitas finita,
absurda autem, si inter illos limites $y$ alicubi infinita evadit. Si integrale tale
$\int\eta\,\mathrm d\xi$, quod generaliter loquendo exhibet aream inter lineam abscissarum atque
curvam, cuius ordinata $=\eta$ pro abscissa $\xi$, secundum regulas suetas evolvimus,
continuitatis immemores, saepissime contradictionibus implicamur. E. g. statuendo
$\eta=\frac1{\xi\xi}$, analysis suppeditat integrale $=C-\frac1\xi$, quo area recte definitur, quamdiu curva continuitatem servat; qua pro $\xi=0$ interrupta, si quis magnitudinem
areae inde ab abscissa negativa usque ad positivam inepte rogat, responsum absurdum a formula feret, eam esse negativam. Quid autem sibi velint haec similiaque
analyseos paradoxa, alia occasione fusius persequemur.

Hic unicam observationem adiicere liceat. Propositis \textit{absque restrictione}
quaestionibus, quae certis casibus absurdae evadere possunt, saepissime ita sibi
consulit analysis, ut responsum ex parte vagum reddat. Ita pro valore integralis
$\iint y\,\mathrm dr\,\mathrm d\varphi$ ab $r=e$ usque ad $r=f$, atque a $\varphi=E$ usque ad $\varphi=F$ extendendi, si valor ipsius $\frac ut$
\[
\begin{aligned}
&\text{pro }r=e,\ \varphi=E\quad\text{designatur per}\quad\theta\\
&\qquad r=e,\ \varphi=F\quad .\quad .\quad .\quad .\quad .\quad .\quad .\quad\theta'\\
&\qquad r=f,\ \varphi=E\quad .\quad .\quad .\quad .\quad .\quad .\quad .\quad\theta''\\
&\qquad r=f,\ \varphi=F\quad .\quad .\quad .\quad .\quad .\quad .\quad .\quad\theta'''
\end{aligned}
\]
per operationes analyticas facile obtinetur
\[
\mathrm{Arc.\ tang}\,\theta-\mathrm{Arc.\ tang}\,\theta'-\mathrm{Arc.\ tang}\,\theta''+\mathrm{Arc.\ tang}\,\theta'''
\]
Revera quidem integrale tunc tantum valorem certum habere potest, quoties $y$
inter limites assignatos semper manet finita: hic valor sub formula tradita utique
contentus, tamen per eam nondum ex asse definitur, quoniam Arc. tang. est functio
multiformis, seorsimque per alias considerationes (haud quidem difficiles) decidere
oportebit, quinam potissimum functionis valores in casu determinato sint adhibendi. Contra quoties $y$ alicubi inter limites assignatos infinita evadit, quaestio
de valore integralis $\iint y\,\mathrm dr\,\mathrm d\varphi$ absurda est: quo non obstante si responsum ab
analysi extorquere obstinaveris, pro methodorum diversitate modo hoc modo illud
reddetur, quae tamen singula sub formula generali ante tradita contenta erunt.
\begin{center}\rule{3cm}{0.3pt}\end{center}
